\documentclass[11pt]{article}

\usepackage[margin=1in]{geometry}
\usepackage{parskip}
\usepackage{hyperref}
\usepackage{times}

\pagestyle{empty}

\begin{document}

\begin{flushright}
\textit{[Date]}  % e.g., February 12, 2026
\end{flushright}

\vspace{0.5em}

\textit{[Editor Name, Degree(s)]}\\
Editor-in-Chief\\
\textit{BMJ Global Health}

\vspace{0.5em}

Dear \textit{[Editor Name]},

We are pleased to submit our manuscript entitled \textbf{``Global, Regional, and National Quality of Care Index for Drug-Susceptible Tuberculosis, 1990--2021: A Systematic Analysis of the Global Burden of Disease Study''} for consideration as an Original Research article in \textit{BMJ Global Health}.

Despite tuberculosis remaining a leading infectious cause of death worldwide, no composite, population-level metric has existed to systematically assess the quality of TB care across countries and over time. Our study addresses this gap by introducing the first Quality of Care Index (QCI) specific to drug-susceptible tuberculosis. Using data from the Global Burden of Disease 2021 Study, we quantify care quality across 204 countries and territories over 31 years (1990--2021) through a novel principal component analysis-based methodology that integrates mortality, morbidity, and disability dimensions into a single interpretable index.

Our findings reveal critical disparities in TB care quality, with the most pronounced gaps in Sub-Saharan Africa, among children under five years of age, and among males. These results are directly relevant to monitoring national and global progress toward the 2030 End TB Strategy targets and the Sustainable Development Goals. The study incorporates rigorous uncertainty quantification and external validation against the Socio-demographic Index, strengthening the robustness and policy relevance of our conclusions.

We believe this work will be of substantial interest to the readership of \textit{BMJ Global Health} given the journal's commitment to advancing health equity and informing global health policy. The QCI framework offers a practical tool for policymakers, programme managers, and researchers seeking to identify where TB care falls short and to prioritise resource allocation accordingly.

This manuscript has not been published previously and is not under consideration by any other journal. All authors have read and approved the submitted version and agree to its submission to \textit{BMJ Global Health}.

We suggest the following potential reviewers with relevant expertise:

\begin{enumerate}
    \item \textit{[Reviewer 1 Name, Degree(s)]}, \textit{[Institution]} --- \textit{[email]}
    \item \textit{[Reviewer 2 Name, Degree(s)]}, \textit{[Institution]} --- \textit{[email]}
    \item \textit{[Reviewer 3 Name, Degree(s)]}, \textit{[Institution]} --- \textit{[email]}
\end{enumerate}

Thank you for your time and consideration. We look forward to your response.

\vspace{0.5em}

Sincerely,

\vspace{1.5em}

\textit{[Corresponding Author Name, Degree(s)]}\\
\textit{[Department / Division]}\\
\textit{[Institution]}\\
\textit{[City, State/Province, Country]}\\
Email: \textit{[email address]}

\end{document}
